\chapter{Glossaire}
\label{chap:glossaire}

\begin{longtable}{p{0.24\linewidth}p{0.68\linewidth}}
\toprule
Terme & Définition \\
\midrule
\endhead
\textbf{DCP} & \emph{DNG Camera Profile}. Profil qui décrit notamment une transformation colorimétrique liée à un boîtier. \\
\textbf{LCP} & \emph{Lens Correction Profile}. Profil de correction optique d'objectif (distorsion, vignetage, aberration selon le profil). \\
\textbf{EV / IL} & Valeur d'exposition. Un incrément de +1 correspond à un doublement de l'énergie lumineuse dans le modèle linéaire. \\
\textbf{HDR} & \emph{High Dynamic Range}. Valeurs de travail au-delà de l'intervalle d'affichage usuel, conservées avant le codage final. \\
\textbf{ICC} & Profil normalisé décrivant la reproduction colorimétrique d'un périphérique ou d'un espace de couleur~\cite{icc2022}. \\
\textbf{LRU} & \emph{Least Recently Used}. Politique d'éviction qui retire l'élément dont le dernier accès est le plus ancien. \\
\textbf{LMS} & Espace de réponses approximatives des cônes Long, Medium et Short, utilisé ici pour l'adaptation chromatique. \\
\textbf{MAE} & \emph{Mean Absolute Error}, moyenne des écarts absolus entre deux images comparées pixel à pixel. \\
\textbf{RAW} & Données issues du capteur, avant le rendu photographique final. \\
\textbf{Sidecar} & Fichier compagnon portant les réglages sans modifier le fichier image source. \\
\textbf{sRGB} & Espace couleur standard destiné notamment aux écrans et au web, défini par IEC~61966-2-1~\cite{srgb}. \\
\textbf{TIFF} & Format conteneur d'image pouvant conserver différentes profondeurs, compressions et profils couleur. \\
\bottomrule
\end{longtable}
